\chapter{СЕРВИСНЫЙ КОНТУР, КОНТЕЙНЕРИЗАЦИЯ И ЭКСПЛУАТАЦИОННАЯ ИНФРАСТРУКТУРА}

\section{gRPC-сервис и формат обмена данными}

Помимо пользовательских приложений, в проекте реализован сервисный контур, предназначенный для интеграции вычислительного ядра во внешнюю инфраструктуру.  
В отличие от desktop-сценариев, здесь основной режим работы связан с пакетной обработкой большого числа точек и включением расчётов в состав других программных систем.  

Сервисный слой реализован отдельным компонентом \texttt{if97\_calculator\_service}.  
Транспортным протоколом выбран gRPC, который обеспечивает строгий контракт на уровне протокола и эффективную сериализацию данных \cite{grpc_docs}.  

Ключевым методом API является двунаправленный потоковый вызов \texttt{If97Service/CalculatePt}.  
Он ориентирован на обработку батчей однотипных запросов, что соответствует природе задачи: множество независимых вычислений $pT$ можно выполнять параллельно и возвращать результаты по мере готовности.  

Формат сообщений построен по схеме Struct-of-Arrays (SoA).  
Смысл SoA состоит в том, что данные передаются не как список объектов «строка таблицы», а как набор примитивных массивов одинаковой длины.  
Например, запрос содержит идентификатор батча и массивы входных параметров:
\begin{itemize}
\item \texttt{PtBatchRequest \{ batch\_id, p[], t[] \}};  
\item \texttt{PtBatchResponse \{ batch\_id, h[], status[], batch\_error\_* \}}.  
\end{itemize}

Преимущества данного подхода:
\begin{itemize}
\item уменьшаются накладные расходы на сериализацию и разбор protobuf-сообщений;  
\item упрощается эффективная data-parallel обработка внутри CPU-пула;  
\item снижается количество промежуточных аллокаций при подготовке данных;  
\item сохраняется возможность построчно возвращать статус (корректно рассчитано/ошибка) и диагностическую информацию.  
\end{itemize}

\section{Модель параллелизма, разделение I/O и CPU, backpressure}

Высокопроизводительная обработка запросов в сервисе требует явного разделения сетевого ввода-вывода и CPU-вычислений.  
Термодинамические расчёты по IF97 являются CPU-bound задачами, поэтому их выполнение в сетевом runtime без изоляции приводит к деградации задержек и пропускной способности.  

В сервисе применена двухконтурная модель исполнения:
\begin{itemize}
\item \texttt{tokio} отвечает за gRPC runtime, сетевые соединения и асинхронный I/O;  
\item \texttt{rayon} выполняет вычисления в CPU-пуле фиксированного размера.  
\end{itemize}

Поток обработки батча включает следующие этапы.  
Сначала сервис принимает очередной фрагмент входного потока, выполняет проверку формата и ограничений.  
Далее батч передаётся в CPU-пул, где вычисления выполняются параллельно по строкам.  
После завершения вычислений результаты возвращаются обратно в поток ответов.  

Для устойчивости под нагрузкой реализован backpressure на уровне количества батчей, одновременно находящихся в обработке.  
Ограничения задаются переменными окружения и используются как предохранители от неконтролируемого роста очередей и потребления памяти:
\begin{itemize}
\item \texttt{IF97\_MAX\_IN\_FLIGHT\_PER\_CONN} --- лимит на одну клиентскую сессию;  
\item \texttt{IF97\_MAX\_IN\_FLIGHT\_GLOBAL} --- глобальный лимит на весь процесс сервиса;  
\item \texttt{IF97\_MAX\_BATCH\_LEN} --- ограничение длины батча;  
\item \texttt{IF97\_GRPC\_MAX\_MSG\_BYTES} --- ограничение размера gRPC-сообщения.  
\end{itemize}

Для эксплуатации также важны health-check и корректное завершение.  
Сервис предоставляет health-статус и поддерживает «drain»-режим при остановке:
перед завершением процесс переводится в состояние \texttt{NotServing}, перестаёт принимать новый трафик и завершает активные вычисления в пределах заданного времени \texttt{IF97\_DRAIN\_SECS}.  

\section{Контейнеризация в Docker}

Для воспроизводимой поставки и упрощения развёртывания сервис упакован в Docker-контейнер \cite{docker_docs}.  
Контейнеризация реализована через многостадийную сборку, в которой компиляция выполняется в отдельной builder-стадии.  
В build-окружение включаются необходимые инструменты (в частности, \texttt{protobuf-compiler} для генерации кода из \texttt{.proto}).  
В итоговый образ попадают только исполняемый файл и минимально необходимое окружение.  

Такой подход даёт практические преимущества:
\begin{itemize}
\item снижает размер конечного образа;  
\item делает поставку независимой от окружения машины;  
\item упрощает запуск сервиса на VPS и в оркестраторах;  
\item обеспечивает единый способ конфигурирования через переменные окружения.  
\end{itemize}

\section{Развёртывание в Kubernetes}

Для эксплуатации в оркестраторе подготовлен набор Kubernetes-манифестов для сервисного слоя \cite{kubernetes_docs}.  
Типовой набор ресурсов включает:
\begin{itemize}
\item \texttt{Deployment} для управления жизненным циклом pod и проведения rolling update;  
\item \texttt{Service} для сетевой публикации gRPC-эндпоинта;  
\item \texttt{HPA} для горизонтального масштабирования по метрикам;  
\item \texttt{PDB} для ограничения деградации доступности при обслуживании узлов;  
\item отдельный \texttt{Job} для нагрузочного тестирования.  
\end{itemize}

При развёртывании важны два эксплуатационных свойства.  
Первое свойство --- масштабируемость: при увеличении нагрузки можно увеличивать число реплик сервиса и настраивать лимиты CPU-пула на уровне конфигурации.  
Второе свойство --- корректная остановка: при обновлениях и перемещениях pod сервис должен завершать активные вычисления предсказуемо, используя health-check и drain.  

\section{CI/CD и автоматизация сборки и релизов}

Автоматизация сборки и поставки реализована через GitHub Actions \cite{github_actions_docs}.  
Цель CI/CD в проекте состоит в том, чтобы гарантировать воспроизводимость артефактов, автоматическую проверку корректности и выпуск релизов под несколько платформ.  

Типовая схема автоматизации включает:
\begin{itemize}
\item workflow валидации: сборка workspace и прогон тестов вычислительного ядра;  
\item workflow публикации документации: сборка rustdoc и размещение документации;  
\item workflow релиза: сборка артефактов под Windows, Linux и macOS и публикация релиза.  
\end{itemize}

Релизные сборки запускаются по тегам формата \texttt{vMAJOR.MINOR.PATCH}, что соответствует требованиям semver.  
Для корректной сборки сервисного слоя в CI учитывается необходимость наличия \texttt{protoc} в Linux runner, поскольку генерация gRPC-обвязки выполняется на этапе сборки.  

В итоговые релизы включаются артефакты пользовательских приложений и сервисной части.  
Это замыкает контур от разработки вычислительного ядра до поставки инструментов, пригодных к установке и эксплуатации.  

\section{Нагрузочное тестирование и эксплуатационные характеристики}

Проверка сервисного контура включает функциональную верификацию и нагрузочные сценарии.  
Функциональная часть опирается на тесты вычислительного ядра \texttt{if97\_core}, поскольку именно оно определяет корректность результатов.  
Нагрузочная часть выполняется через специализированные сценарии (например, \texttt{grpc\_loadtest}) и через Kubernetes Job, что приближает условия проверки к эксплуатационным.  

Ключевыми эксплуатационными характеристиками являются пропускная способность пакетного расчёта и устойчивость под нагрузкой.  
По результатам измерений, для release-сборки достигается throughput порядка $181$ тыс. строк/с при native-запуске и порядка $148$ тыс. строк/с при запуске в Docker-контейнере.  
Снижение пропускной способности в контейнере отражает допустимый инфраструктурный оверхед при сохранении высокой производительности вычислительного ядра.  

Устойчивость сервиса обеспечивается сочетанием инженерных решений:
разделением I/O и CPU, использованием SoA-формата данных, ограничениями in-flight батчей, контролем размеров сообщений и корректным жизненным циклом (health-check и drain).  
